\section{Legal Domain}
\label{sec:domains}
In this section, we describe the Italian legal system, the different types of trials with the documents involved, and the relationship between Italian law, \acl{ai} and predictability (ai act, gdpr).
italian criminal proceedings, vs civil
\todo{introdurre ai prima}
\todo{different law codes, civil law, supreme court (vs common law)}

The Italian Republic uses the \emph{civil law} system, in contrast with countries such as the United Kingdom of Great Britain and Northern Ireland and the United State of America, using the \emph{common law} system. Civil law is based on codes, while common law is based on case-law~\cite{Carlson2009}. Typically, civil law systems rely on four foundational codes: civil code, the civil procedure code, criminal code, and criminal procedure code. \todo{here add a brief descriptions of what are these codes} These codes are intended as coherent and consistent principles for supporting the broader legal system. In contrast, common law systems typically develop codes in response to specific issues, often drafting them only after principles have already emerged through case law~\cite{Carlson2009}.
Furthermore, in common law, precedent cases are of central importance, as judges are required to follow earlier decisions, and a single ruling from a higher court may determine the outcome of a case. This is different in civil law, where a single precedent case is generally not treated as binding, with the possible exception of decisions in the constitutional sphere~\cite{Carlson2009}. 

The organizational structure of civil law court systems is broadly similar to that of common law systems, featuring a court of first instance, a court of appeals, and a supreme court. Additionally, specialized courts exist with jurisdiction limited to particular areas of law. Unlike in common law jurisdictions, juries are generally not used in civil law systems. The cases are decided by judges, eventually including lay judges~\cite{Carlson2009}, such as in the Italian court of assize~\cite[Article 3]{legge10apr1951}---which is competent to judge serious criminal offenses, e.g., the ones carrying maximum penalties of at least 24 years of imprisonment~\cite[Article 5]{CPP}.  

In the Italian system, both in civil and in criminal procedures, the three grades of jurisdiction operate as follows. The courts of first instance decide the facts and merits of a case in their entirety~\cite[Articles 99--310]{CPC}; \cite[Articles 429--544]{CPP}. The courts of appeal then provide a second, comprehensive review to ensure both factual and legal correctness~\cite[Articles 339--359]{CPC}; \cite[Articles 593--605]{CPP}, thereby giving effect to the constitutional guarantee of challenging judicial decisions~\cite[Article 111]{Costituzione}. Finally, the Court of Cassation represents the third and highest grade, serving as the supreme authority tasked not with re-examining facts but with ensuring the uniform interpretation and application of the law, as regulated in the codes of procedure~\cite[Articles 360--382]{CPC}; \cite[Articles 606--613]{CPP}.
Decisions of the Court of Cassation are not laws and are not formally binding, yet they strongly influence lower courts and promote a consistent interpretation of the law, reflecting the court's so-called \emph{nomophilactic function}~\cite{lazovic2022}. In Italy, the supreme court has a dual role: to render a final judgment on individual cases, thereby ensuring justice, and to establish interpretative guidelines that shape national jurisprudence~\cite[Article 65]{RD1941_12_art65}.


\subsection{Legal Proceedings in the Italian Legal System}

In the Italian legal system, a legal proceeding (\textit{procedimento}) may arise in either the civil or the criminal sphere, with distinct initiation mechanisms reflecting the nature of the dispute.

% \subsubsection{Civil Proceedings}

Civil proceedings are initiated by a private party. The \emph{plaintiff} begins the process by filing an writ of summons with the competent court. This document must indicate: the identities of the parties, the tribunal being addressed, the factual and legal basis of the claim, the evidence intended to be produced, and the date of the hearing \cite[Article 163]{CPC}. 

The \emph{defendant} responds with a statement of defense, outlining their defenses and any counterclaims \cite[Article 167]{CPC}. From this moment, the judicial authority schedules hearings, manages the collection of evidence, and ultimately renders a decision---a \emph{judgment}---which must include motivations as required by the Constitution~\cite[Article 111]{Costituzione}.

% \subsubsection{Criminal Proceedings}

Criminal proceedings begin with the occurrence of a potential criminal offense. When the report of a crime is received, the public prosecutor opens the preliminary investigation~\cite[Article 326]{CPP}.

During this phase, investigative acts are performed by the public prosecutor and the judicial police, and documented in official records~\cite[Article 373]{CPP}. The suspect is formally notified through a guarantee notic, a document informing them of the charges and their rights~\cite[Article 369]{CPP}. If sufficient evidence exists, the prosecutor requests an indictment and the judge for the preliminary hearing decides whether the case proceeds to trial \cite[Article 416]{CPP}.

\todo{ai and legal, plus predictability}
\todo{per i documenti: sentenze e fascicoli. vedi chat e tesi di alessandro mariani. 132 cpc}
\todo{corti europee e internazionali e impatti sul legal ita: gdpr ai act anche}

% \subsubsection{Documents}
The documentation of civil and criminal proceedings is governed by the provisions of the respective codes of procedure~\cite[Articles 168--169]{CPC}; \cite[Articles 431, 114, 391-octies]{CPP}, which govern the formation, content, and accessibility of case files.
Civil files typically include the summons, payment receipts, pleadings, briefs, hearing minutes, judicial orders, evidentiary acts, and judgments~\cite[Article 168]{CPC}, while criminal trial files contain acts related to prosecutability and civil claims, non-repeatable acts by police or parties, evidentiary acts, official documents, and, when applicable, seized items~\cite[Article 431]{CPP}.
Furthermore, the content of the judgment are regulated, requiring identification of the parties, a statement of facts, the legal reasoning, and the operative part of the decision~\cite[Articles 544--545]{CPP}; \cite[Article 133]{CPC}. For this reason, judgments can be treated as semi-structured documents, in the sense that they usually contain identifiable sections, even though the expression and organization of legal language may vary---e.g., depending of the style of the judge~\cite{juliahb2024,galli2025analisi,pont2023legalsummarisationllmsprodigit}.

\todo{international influence, eu gdpr ai act}
\todo{definition of AI}
\todo{why applying ai to legal, maybe focus what we did}

% \subsection{Analyzing Legal Documents}
% \todo{also AI for legal}

% improving access to legal information, facilitating search and retrieval of relevant case law, generating summaries and keywords, extracting arguments, and even offering statistical forecasts of outcomes. These functions promise to enhance efficiency in overloaded court systems, reduce inconsistencies in jurisprudence, and support judges and lawyers in handling the growing volume and complexity of legal data. Importantly, AI is conceived not as a substitute for judicial decision-making, but as an auxiliary instrument that can strengthen transparency, consistency, and access to justice~\cite{juliahb2024}.
% ADELE project and outcome prediction
\todo{distinguish between "retrieval/information tools" and decision systems}
\todo{also speak about decision systems in healthcare}.

% chat
% **Why AI can be useful in justice**
% Artificial intelligence (AI) is increasingly seen as a valuable support tool for judicial systems. According to the JuLIA Handbook, AI can assist in a variety of tasks: improving access to legal information, facilitating search and retrieval of relevant case law, generating summaries and keywords, extracting arguments, and even offering statistical forecasts of outcomes. These functions promise to enhance efficiency in overloaded court systems, reduce inconsistencies in jurisprudence, and support judges and lawyers in handling the growing volume and complexity of legal data. Importantly, AI is conceived not as a substitute for judicial decision-making, but as an auxiliary instrument that can strengthen transparency, consistency, and access to justice.

% **What has been done so far in Italy**
% Italy has begun experimenting with AI tools, especially in the context of predictive justice. The Documentation Centre (CED) of the Court of Cassation has explored algorithmic methods for citation analysis and legal research, while research projects such as ADELE have tested applications for summarization, argument extraction, and outcome prediction. The 2023 Cartabia Reform introduced a shift from “documents” to “data,” laying the groundwork for more systematic use of AI in judicial administration. Legislative attention has also increased: Italy presented a proposal (Bill no. 1146/2024) to regulate AI specifically in the justice sector, aligning with the EU Artificial Intelligence Act but adopting a cautious approach that limits AI’s role to supporting research and organizational tasks.

% **Issues and regulatory safeguards**
% The JuLIA Handbook underscores that AI in justice is a high-risk domain, requiring strong safeguards. Risks include algorithmic opacity, biased datasets, and potential erosion of judicial independence if predictive tools were given too much weight. For this reason, both European and Italian frameworks emphasize principles of transparency, non-discrimination, due process, and human oversight. The Italian Council of State has further clarified that algorithmic systems must comply with constitutional requirements of impartiality and accountability. Overall, Italy’s approach reflects a careful balance: acknowledging AI’s potential to make justice more efficient and accessible, while embedding safeguards to protect judicial autonomy and fundamental rights.
%%%

\todo{add citations and articles}
\todo{later expand this with more paper on automated decisions, legal simulations (see e-mails),... but maybe extraction papers and entity-based later}
\todo{ieri approvata legge ai in italia!!!}
% https://www.wired.it/article/ddl-ai-senato-approvazione-regole-controllo-autorita/
% https://www.repubblica.it/tecnologia/2025/09/18/news/legge_italiana_intelligenza_artificiale_deep_fake-424854661/
\subsection{Artificial Intelligence in the Legal Domain}

\ac{ai} offers considerable opportunities for the legal sector, fundamentally reshaping traditional workflows through the automation of routine tasks and the enhancement of complex decision-making processes \cite{juliahb2024}. The application of \ac{ai} is not intended to supplant human legal expertise but rather to act as a sophisticated form of decision-support systems (DSSs), augmenting the capabilities of legal professionals. Such systems can perform an array of tasks, from conducting advanced precedent searches to predicting potential case outcomes, thereby improving efficiency while maintaining the integrity of human judicial oversight \cite{juliahb2024}.

A notable example of a practical \ac{ai} tool is the Analytics for Deciding Legal Cases (ADELE) project~\cite{galli2024analytics}, which focused on a variety of legal domains, including value-added tax (VAT), and trademarks and patents (T\&P) cases. The project's objectives centered on supporting legal research and decision-making through a suite of functionalities. These included the automatic extraction of citations and the subsequent creation of citation networks, which aids in identifying influential cases and clusters of similar jurisprudence. Furthermore, ADELE utilized machine learning models for the extraction of summaries and keywords from legal texts, thereby streamlining information retrieval and organizational processes.\todo{more focus on the utility of speeding up justice} A particularly ambitious aspect of the project was the experimentation with argument extraction and the predictive analysis of case outcomes, which sought to establish correlations between the arguments presented by litigants and the final judgments \cite{juliahb2024}.

In the context of information retrieval and organization, the importance of data annotation cannot be overstated. Projects in Italy, such as the PNRR-PON ``Giustizia Agile'' (Italian for ``agile justice''), have focused on creating gold standards for Italian legal documents to support further \ac{ai} research. This included the crucial task of Named Entity Recognition (NER), which is foundational for enabling \ac{ai} systems to accurately identify and classify key entities within legal texts, such as names of persons, organizations, or legal provisions, thereby improving the precision of retrieval mechanisms \cite{juliahb2024}.

\subsection{\ac{ai} Regulations}

The integration of artificial intelligence (\ac{ai}) into the legal and judicial system presents significant challenges, particularly concerning fundamental rights and ethical principles. The European Union has addressed these concerns through the Artificial Intelligence Act (\ac{ai} Act)~\cite{ai_act}, which classifies \ac{ai} systems used in the administration of justice as \say{high-risk.} While the \ac{ai} Act does not prohibit their use, it requires that the final decision-making authority remains with a human, ensuring that \ac{ai} functions as a support mechanism rather than a final arbiter \cite{juliahb2024}. This aligns with the European Ethical Charter~\cite{ethicalcharter2018} on the use of \ac{ai} in judicial systems, which emphasizes principles such as transparency, non-discrimination, and due process. The principle of transparency is particularly critical, as it necessitates that the reasoning behind an algorithm's output can be understood and explained to justify a judicial decision.

These principles have been further consolidated in Italy with the approval of the Law on Artificial Intelligence (L. 1146/2024)~\cite{leggeia2025}. This national law integrates and strengthens the European regulatory framework, establishing Italy as one of the first EU countries to have specific legislation on the matter~\cite{reuters_italy_ai_law_2025}. The Italian law reaffirms the anthropocentric approach, ensuring that \ac{ai} is used in the judicial sector exclusively for instrumental and support activities. Furthermore, it introduces important criminal provisions, such as penalties for the illicit dissemination of content generated or manipulated with \ac{ai}, thereby underscoring the law's role in mitigating emerging risks.

These challenges and regulations have already been reflected in case law across both Italy and the EU. The Tribunal of Bologna, for example, has considered cases involving algorithmic decision-making, basing its analysis on the principles of non-discrimination and equality. Similarly, the Court of Justice of the European Union (CJEU) in the SCHUFA Holding case~\cite{schufaholding2023} examined whether an \ac{ai}-driven credit scoring system constitutes a form of automated decision-making under the GDPR (General Data Protection Regulation)~\cite{gdpr2016}, a precedent with significant implications for the use of \ac{ai} in judicial contexts \cite{juliahb2024}. These cases demonstrate the ongoing legal and ethical discourse surrounding \ac{ai} and the urgent need for robust regulatory frameworks to govern its application in the justice system.\todo{da verificare. scritto da gemini}
\todo{aggiungi altri progetti legali tipo datalake e magari cose su extrazione. o dopo?}

\todo{rw from aixia,wise,claw}


\subsection{Judgments}
Legal documents---particularly court judgments and orders---hold information that can serve many applications (such as legal case retrieval and process analysis), fulfill different purposes (ranging from comparing similar cases to promote consistent rulings, to identifying trends within specific legal domains, e.g., assessing average maintenance payments in divorces relative to the spouses’ financial conditions), and support a wide range of stakeholders (including judges, court administrators, lawmakers, legal practitioners, researchers, and the public)~\cite{BELLANDI2024}\todo{check}.
An example of fictional Italian civil judgment, translated in English, is available in \Cref{fig:civil-judgment-legal-articles}.

\begin{figure}[htbp]
\centering
\begin{minipage}{0.9\linewidth}
\small
ORDINARY COURT OF TURIN\\
Civil Division\\
Judgment No. 1234/2025\\
Published on March 15, 2025\\[1em]
ITALIAN REPUBLIC\\
IN THE NAME OF THE ITALIAN PEOPLE\\[1em]
The Court, sitting as a single judge, in the person of Judge Andrea Bianchi, has delivered the following:\\[1em]
JUDGMENT\\[0.5em]
in the civil case registered under No. 4567/2024\\
between\\[0.5em]
Mario Rossi, residing in Turin, represented and defended by Attorney Lucia Verdi, with elected domicile at her office in Turin -- plaintiff\\[0.5em]
and\\[0.5em]
Carlo Neri, residing in Milan, represented and defended by Attorney Marco Ferrari, with elected domicile at his office in Milan -- defendant\\[1em]
FACTS\\
The plaintiff summoned the defendant to court, alleging that on January 10, 2023, he lent him the sum of EUR 15,000, as evidenced by a private written agreement signed by both parties, with the obligation to repay by June 30, 2023, pursuant to Article 1813 c.c. (loan contract). Despite repeated requests, the defendant allegedly failed to repay the sum. The defendant, appearing in court at the Court Building in Turin, contested the existence of the debt, claiming that the sum received was a gift under Article 769 c.c.\\[1em]
REASONS FOR THE DECISION\\
The evidence confirmed the existence of a loan contract under Article 1813 c.c., as established by the private written agreement submitted in evidence and by the testimony of two independent witnesses (pursuant to Articles 2721--2725 c.c. on testimonial evidence). The claim of a gift is unfounded, as no donative intent on the part of the plaintiff was proven, as required by Article 769 of the Italian Civil Code. Pursuant to Article 2697 c.c., the burden of proof lies with the party asserting release from an obligation, and in this case the defendant failed to provide adequate proof. Therefore, the plaintiff's credit must be recognized in the amount of EUR 15,000, plus statutory interest from July 1, 2023, in accordance with Article 1282 c.c.\\[1em]
RULING\\
The Court, definitively ruling, dismissing any other claims or objections,\\
orders Mr. Carlo Neri to pay Mr. Mario Rossi the sum of EUR 15,000, plus statutory interest from July 1, 2023 until payment in full (Articles 1282 and 1284 c.c.);\\
orders the defendant to reimburse the plaintiff's legal costs, set at EUR 2,500 for legal fees, plus statutory surcharges, pursuant to Articles 91 and 92 c.p.c.\\[1em]
Turin, March 10, 2025\\
The Judge\\
Andrea Bianchi
\end{minipage}
\caption{Illustrative example of an Italian civil judgment. c.c. stands for Italian Civil Code~\cite{CC} and c.p.c. for Italian Code of Civil Procedure~\cite{CPC}.}
\label{fig:civil-judgment-legal-articles}
\end{figure}



\todo{something general for legal}

\cite{BELLANDI2024,pozzi_wise_2025,aixia_2023,Agazzi2025DAVE,real_batini2021semantic}
\todo{italian legal system. civile penale, investigations. usa tesi dottorato legal}
\todo{AI act. news for legal AI italy, batini,...}
\subsection{Criminal Investigations}
Data in criminal investigations comes from a heterogeneity of sources such as documents, e.g., reports by consultants, (semi-)structured files, e.g., phone call records, bank transfers, and digitalized papers, e.g., seized in the investigation~\cite{pozzi_wise_2025}.\todo{check}

\todo{focus on IMA}

\begin{figure}[htbp]
\centering
\begin{minipage}{0.92\linewidth}
\small
\textbf{Chat Transcript (Extract)}\\
Participants: Mario Rossi, Luca Bianchi\\
Date Range: 2025-07-28 to 2025-07-30\\[0.5em]
\textbf{[2025-07-28 -- 14:12]}  Mario Rossi: Hey, are we still meeting on Friday?\\
\textbf{[2025-07-28 -- 14:13]} Luca Bianchi: Yeah, that's August 1st, right?\\
\textbf{[2025-07-28 -- 14:14]}  Mario Rossi: Correct.\\
\textbf{[2025-07-28 -- 14:15]} Luca Bianchi: Where though? The place near Piazza Vittorio Veneto in Turin or the one by the river?\\
\textbf{[2025-07-28 -- 14:17]}  Mario Rossi: Let's do the bar at Via Roma 210, easier to park.\\
\textbf{[2025-07-28 -- 14:18]} Luca Bianchi: Fine. Anyone else coming?\\
\textbf{[2025-07-28 -- 14:19]}  Mario Rossi: Maybe Gianni, but he's in Milan until Thursday.\\[0.5em]
\textbf{[2025-07-29 -- 08:45]}  Mario Rossi: \textit{[Voice message -- 2 min 14 sec]}\\
\quad \textit{Transcription:}\\
\quad So listen, I was thinking\dots it might be better if we meet a bit earlier, like around 6:30 PM, because I have to talk to you about something before Gianni gets there. It's about that thing we discussed last time at Marco's place — you know, the one near Porta Susa station. I also want to avoid too many people seeing us together at the main bar during the busy hour. If we meet earlier, we can grab a table at the back, order something small, and then, when Gianni arrives, we just act like we've been casually catching up. Also, there's a parking spot I found behind the cinema on Via Po, usually empty at that time. Oh, and bring the envelope I gave you last month, I'll need to check something inside. Don't text me too much about this, we'll talk in person.\\[0.5em]
\textbf{[2025-07-29 -- 09:02]} Luca Bianchi: Got it. 6:30 PM, Via Roma 210, Friday. I'll bring the envelope.\\
\textbf{[2025-07-30 -- 16:47]}  Mario Rossi: Perfect.
\end{minipage}
\caption{Illustrative example of chat exchange.}
\label{fig:chat-investigation-example}
\end{figure}


\todo{guarda 23 di claw survey ai in legal. also ner on turkish [3]}